
\documentclass{article}
\usepackage{colortbl}
\usepackage{makecell}
\usepackage{multirow}
\usepackage{supertabular}

\begin{document}

\newcounter{utterance}

\twocolumn

{ \footnotesize  \setcounter{utterance}{1}
\setlength{\tabcolsep}{0pt}
\begin{supertabular}{c@{$\;$}|p{.15\linewidth}@{}p{.15\linewidth}p{.15\linewidth}p{.15\linewidth}p{.15\linewidth}p{.15\linewidth}}

    \# & $\;$A & \multicolumn{4}{c}{Game Master} & $\;\:$B\\
    \hline 

    \theutterance \stepcounter{utterance}  

    & & \multicolumn{4}{p{0.6\linewidth}}{\cellcolor[rgb]{0.9,0.9,0.9}{%
	\makecell[{{p{\linewidth}}}]{% 
	  \tt {\tiny [A$\langle$GM]}  
	 State giocando a un gioco di negoziazione in cui dovete accordarvi su come dividere una serie di oggetti tra voi e un altro giocatore.\\ \tt \\ \tt Regole:\\ \tt (a) A Lei e all'altro giocatore viene dato un insieme di oggetti. Ognuno di voi riceve anche una funzione di valore segreta, che rappresenta il valore di ciascun tipo di oggetto.\\ \tt (b) Si scambiano messaggi con l'altro giocatore per concordare chi si aggiudica gli oggetti. Potete inviare un massimo di 5 messaggi ciascuno, oppure terminare in anticipo facendo una proposta segreta in qualsiasi momento.\\ \tt (c) A ciascuno di voi viene chiesto di inviare una proposta segreta indicando gli oggetti che desiderate, formattata tra parentesi quadre come segue: "[Proposta: <numero> <nome oggetto>, <numero> <nome oggetto>, <...>]".\\ \tt (d) Se le vostre proposte sono complementari, cioè ci sono abbastanza oggetti per soddisfare entrambe le proposte, a ciascun giocatore viene assegnato un punteggio basato sulla somma dei valori degli oggetti ricevuti. In caso contrario, entrambi ricevono zero punti.\\ \tt (e) L'obiettivo è massimizzare il proprio punteggio.\\ \tt \\ \tt Cominciamo.\\ \tt \\ \tt L'insieme degli oggetti disponibili è:\\ \tt \\ \tt 1 lampada, 3 ditali, 1 corda\\ \tt \\ \tt La funzione valore segreta è:\\ \tt \\ \tt lampada: 3, ditale: 2, corda: 1\\ \tt \\ \tt IMPORTANTE: i vostri messaggi, a meno che non si tratti di una proposta segreta, vengono trasmessi direttamente all'altro giocatore, quindi non includete alcuna risposta alle regole o testo di annuncio del vostro messaggio. Per fare una proposta segreta, utilizzate il formato indicato. Non utilizzare delle parentesi quadre quando si comunica all'altro giocatore, altrimenti verrà interpretata come una proposta segreta.\\ \tt  
	  } 
	   } 
	   } 
	 & \\ 
 

    \theutterance \stepcounter{utterance}  

    & \multicolumn{4}{p{0.6\linewidth}}{\cellcolor[rgb]{0.8,1,0.9}{%
	\makecell[{{p{\linewidth}}}]{% 
	  \tt {\tiny [A$\rangle$GM]}  
	 Ciao, sono pronto a iniziare. Vorrei sapere cosa pensi di dividere gli oggetti. Io penso che potremmo iniziare dividendo i ditali, visto che ce ne sono 3. Che ne pensi? 
	  } 
	   } 
	   } 
	 & & \\ 
 

    \theutterance \stepcounter{utterance}  

    & & & \multicolumn{2}{p{0.3\linewidth}}{\cellcolor[rgb]{0.95,0.95,0.95}{%
	\makecell[{{p{\linewidth}}}]{% 
	  \tt {\tiny [GM$|$GM]}  
	 continue 
	  } 
	   } 
	   } 
	 & & \\ 
 

    \theutterance \stepcounter{utterance}  

    & & & \multicolumn{2}{p{0.3\linewidth}}{\cellcolor[rgb]{0.95,0.95,0.95}{%
	\makecell[{{p{\linewidth}}}]{% 
	  \tt {\tiny [GM$|$GM]}  
	 Ciao, sono pronto a iniziare. Vorrei sapere cosa pensi di dividere gli oggetti. Io penso che potremmo iniziare dividendo i ditali, visto che ce ne sono 3. Che ne pensi? 
	  } 
	   } 
	   } 
	 & & \\ 
 

    \theutterance \stepcounter{utterance}  

    & & \multicolumn{4}{p{0.6\linewidth}}{\cellcolor[rgb]{0.9,0.9,0.9}{%
	\makecell[{{p{\linewidth}}}]{% 
	  \tt {\tiny [GM$\rangle$B]}  
	 State giocando a un gioco di negoziazione in cui dovete accordarvi su come dividere una serie di oggetti tra voi e un altro giocatore.\\ \tt \\ \tt Regole:\\ \tt (a) A Lei e all'altro giocatore viene dato un insieme di oggetti. Ognuno di voi riceve anche una funzione di valore segreta, che rappresenta il valore di ciascun tipo di oggetto.\\ \tt (b) Si scambiano messaggi con l'altro giocatore per concordare chi si aggiudica gli oggetti. Potete inviare un massimo di 5 messaggi ciascuno, oppure terminare in anticipo facendo una proposta segreta in qualsiasi momento.\\ \tt (c) A ciascuno di voi viene chiesto di inviare una proposta segreta indicando gli oggetti che desiderate, formattata tra parentesi quadre come segue: "[Proposta: <numero> <nome oggetto>, <numero> <nome oggetto>, <...>]".\\ \tt (d) Se le vostre proposte sono complementari, cioè ci sono abbastanza oggetti per soddisfare entrambe le proposte, a ciascun giocatore viene assegnato un punteggio basato sulla somma dei valori degli oggetti ricevuti. In caso contrario, entrambi ricevono zero punti.\\ \tt (e) L'obiettivo è massimizzare il proprio punteggio.\\ \tt \\ \tt Cominciamo.\\ \tt \\ \tt L'insieme degli oggetti disponibili è:\\ \tt \\ \tt 1 lampada, 3 ditali, 1 corda\\ \tt \\ \tt La funzione valore segreta è:\\ \tt \\ \tt lampada: 3, ditale: 1, corda: 4\\ \tt \\ \tt IMPORTANTE: i vostri messaggi, a meno che non si tratti di una proposta segreta, vengono trasmessi direttamente all'altro giocatore, quindi non includete alcuna risposta alle regole o testo di annuncio del vostro messaggio. Per fare una proposta segreta, utilizzate il formato indicato. Non utilizzare delle parentesi quadre quando si comunica all'altro giocatore, altrimenti verrà interpretata come una proposta segreta.\\ \tt \\ \tt \\ \tt Ciao, sono pronto a iniziare. Vorrei sapere cosa pensi di dividere gli oggetti. Io penso che potremmo iniziare dividendo i ditali, visto che ce ne sono 3. Che ne pensi? 
	  } 
	   } 
	   } 
	 & \\ 
 

    \theutterance \stepcounter{utterance}  

    & & & \multicolumn{4}{p{0.6\linewidth}}{\cellcolor[rgb]{1,0.85,0.72}{%
	\makecell[{{p{\linewidth}}}]{% 
	  \tt {\tiny [GM$\langle$B]}  
	 Sembra una buona idea iniziare con i ditali, visto che sono il tipo di oggetto più numeroso. Potremmo dividerli equamente, quindi ognuno di noi potrebbe prendere un ditale e un altro oggetto, o anche due ditali e mettere da parte gli altri oggetti per discuterli in seguito. Cosa ne pensi di fare con la lampada e la corda? 
	  } 
	   } 
	   } 
	  \\ 
 

    \theutterance \stepcounter{utterance}  

    & & & \multicolumn{2}{p{0.3\linewidth}}{\cellcolor[rgb]{0.95,0.95,0.95}{%
	\makecell[{{p{\linewidth}}}]{% 
	  \tt {\tiny [GM$|$GM]}  
	 continue 
	  } 
	   } 
	   } 
	 & & \\ 
 

    \theutterance \stepcounter{utterance}  

    & & & \multicolumn{2}{p{0.3\linewidth}}{\cellcolor[rgb]{0.95,0.95,0.95}{%
	\makecell[{{p{\linewidth}}}]{% 
	  \tt {\tiny [GM$|$GM]}  
	 Sembra una buona idea iniziare con i ditali, visto che sono il tipo di oggetto più numeroso. Potremmo dividerli equamente, quindi ognuno di noi potrebbe prendere un ditale e un altro oggetto, o anche due ditali e mettere da parte gli altri oggetti per discuterli in seguito. Cosa ne pensi di fare con la lampada e la corda? 
	  } 
	   } 
	   } 
	 & & \\ 
 

    \theutterance \stepcounter{utterance}  

    & & \multicolumn{4}{p{0.6\linewidth}}{\cellcolor[rgb]{0.9,0.9,0.9}{%
	\makecell[{{p{\linewidth}}}]{% 
	  \tt {\tiny [A$\langle$GM]}  
	 Sembra una buona idea iniziare con i ditali, visto che sono il tipo di oggetto più numeroso. Potremmo dividerli equamente, quindi ognuno di noi potrebbe prendere un ditale e un altro oggetto, o anche due ditali e mettere da parte gli altri oggetti per discuterli in seguito. Cosa ne pensi di fare con la lampada e la corda? 
	  } 
	   } 
	   } 
	 & \\ 
 

    \theutterance \stepcounter{utterance}  

    & \multicolumn{4}{p{0.6\linewidth}}{\cellcolor[rgb]{0.8,1,0.9}{%
	\makecell[{{p{\linewidth}}}]{% 
	  \tt {\tiny [A$\rangle$GM]}  
	 La tua idea di dividere i ditali equamente mi sembra giusta. Per quanto riguarda la lampada e la corda, credo che la lampada sia un oggetto più prezioso e quindi potrebbe essere il pezzo chiave da aggiudicare. La corda sembra meno importante, ma potrebbe comunque essere utile. Se dovessi scegliere, prenderei la lampada e ti lascerei la corda, ma sono aperto a negoziare. Che ne pensi se io prendessi la lampada e due ditali, e tu prendessi la corda e l'altro ditale? 
	  } 
	   } 
	   } 
	 & & \\ 
 

    \theutterance \stepcounter{utterance}  

    & & & \multicolumn{2}{p{0.3\linewidth}}{\cellcolor[rgb]{0.95,0.95,0.95}{%
	\makecell[{{p{\linewidth}}}]{% 
	  \tt {\tiny [GM$|$GM]}  
	 continue 
	  } 
	   } 
	   } 
	 & & \\ 
 

    \theutterance \stepcounter{utterance}  

    & & & \multicolumn{2}{p{0.3\linewidth}}{\cellcolor[rgb]{0.95,0.95,0.95}{%
	\makecell[{{p{\linewidth}}}]{% 
	  \tt {\tiny [GM$|$GM]}  
	 La tua idea di dividere i ditali equamente mi sembra giusta. Per quanto riguarda la lampada e la corda, credo che la lampada sia un oggetto più prezioso e quindi potrebbe essere il pezzo chiave da aggiudicare. La corda sembra meno importante, ma potrebbe comunque essere utile. Se dovessi scegliere, prenderei la lampada e ti lascerei la corda, ma sono aperto a negoziare. Che ne pensi se io prendessi la lampada e due ditali, e tu prendessi la corda e l'altro ditale? 
	  } 
	   } 
	   } 
	 & & \\ 
 

    \theutterance \stepcounter{utterance}  

    & & \multicolumn{4}{p{0.6\linewidth}}{\cellcolor[rgb]{0.9,0.9,0.9}{%
	\makecell[{{p{\linewidth}}}]{% 
	  \tt {\tiny [GM$\rangle$B]}  
	 La tua idea di dividere i ditali equamente mi sembra giusta. Per quanto riguarda la lampada e la corda, credo che la lampada sia un oggetto più prezioso e quindi potrebbe essere il pezzo chiave da aggiudicare. La corda sembra meno importante, ma potrebbe comunque essere utile. Se dovessi scegliere, prenderei la lampada e ti lascerei la corda, ma sono aperto a negoziare. Che ne pensi se io prendessi la lampada e due ditali, e tu prendessi la corda e l'altro ditale? 
	  } 
	   } 
	   } 
	 & \\ 
 

    \theutterance \stepcounter{utterance}  

    & & & \multicolumn{4}{p{0.6\linewidth}}{\cellcolor[rgb]{1,0.85,0.72}{%
	\makecell[{{p{\linewidth}}}]{% 
	  \tt {\tiny [GM$\langle$B]}  
	 La tua proposta non sembra male, ma non sono sicuro che sia la scelta migliore per me. Prendere la corda e solo un ditale potrebbe non essere abbastanza, considerando che tu avresti la lampada e due ditali. Vorrei provare a trovare un equilibrio migliore. Che ne pensi se invece io prendessi la lampada e un ditale, e tu prendessi la corda e due ditali? In questo modo, entrambi avremmo un oggetto più "prezioso" e una quantità equa di ditali. 
	  } 
	   } 
	   } 
	  \\ 
 

    \theutterance \stepcounter{utterance}  

    & & & \multicolumn{2}{p{0.3\linewidth}}{\cellcolor[rgb]{0.95,0.95,0.95}{%
	\makecell[{{p{\linewidth}}}]{% 
	  \tt {\tiny [GM$|$GM]}  
	 continue 
	  } 
	   } 
	   } 
	 & & \\ 
 

    \theutterance \stepcounter{utterance}  

    & & & \multicolumn{2}{p{0.3\linewidth}}{\cellcolor[rgb]{0.95,0.95,0.95}{%
	\makecell[{{p{\linewidth}}}]{% 
	  \tt {\tiny [GM$|$GM]}  
	 La tua proposta non sembra male, ma non sono sicuro che sia la scelta migliore per me. Prendere la corda e solo un ditale potrebbe non essere abbastanza, considerando che tu avresti la lampada e due ditali. Vorrei provare a trovare un equilibrio migliore. Che ne pensi se invece io prendessi la lampada e un ditale, e tu prendessi la corda e due ditali? In questo modo, entrambi avremmo un oggetto più "prezioso" e una quantità equa di ditali. 
	  } 
	   } 
	   } 
	 & & \\ 
 

    \theutterance \stepcounter{utterance}  

    & & \multicolumn{4}{p{0.6\linewidth}}{\cellcolor[rgb]{0.9,0.9,0.9}{%
	\makecell[{{p{\linewidth}}}]{% 
	  \tt {\tiny [A$\langle$GM]}  
	 La tua proposta non sembra male, ma non sono sicuro che sia la scelta migliore per me. Prendere la corda e solo un ditale potrebbe non essere abbastanza, considerando che tu avresti la lampada e due ditali. Vorrei provare a trovare un equilibrio migliore. Che ne pensi se invece io prendessi la lampada e un ditale, e tu prendessi la corda e due ditali? In questo modo, entrambi avremmo un oggetto più "prezioso" e una quantità equa di ditali. 
	  } 
	   } 
	   } 
	 & \\ 
 

    \theutterance \stepcounter{utterance}  

    & \multicolumn{4}{p{0.6\linewidth}}{\cellcolor[rgb]{0.8,1,0.9}{%
	\makecell[{{p{\linewidth}}}]{% 
	  \tt {\tiny [A$\rangle$GM]}  
	 La tua controproposta è interessante e sembra più equilibrata. Io avrei la corda e due ditali, mentre tu avresti la lampada e un ditale. Questo modo di dividere gli oggetti sembra più giusto e potrebbe funzionare. Sono disposto ad accettare questa divisione, ma prima di prendere una decisione definitiva, vorrei chiederti: sei soddisfatto di questa proposta o ci sono ancora delle possibili modifiche che potremmo apportare per rendere la divisione ancora più equa? 
	  } 
	   } 
	   } 
	 & & \\ 
 

    \theutterance \stepcounter{utterance}  

    & & & \multicolumn{2}{p{0.3\linewidth}}{\cellcolor[rgb]{0.95,0.95,0.95}{%
	\makecell[{{p{\linewidth}}}]{% 
	  \tt {\tiny [GM$|$GM]}  
	 continue 
	  } 
	   } 
	   } 
	 & & \\ 
 

    \theutterance \stepcounter{utterance}  

    & & & \multicolumn{2}{p{0.3\linewidth}}{\cellcolor[rgb]{0.95,0.95,0.95}{%
	\makecell[{{p{\linewidth}}}]{% 
	  \tt {\tiny [GM$|$GM]}  
	 La tua controproposta è interessante e sembra più equilibrata. Io avrei la corda e due ditali, mentre tu avresti la lampada e un ditale. Questo modo di dividere gli oggetti sembra più giusto e potrebbe funzionare. Sono disposto ad accettare questa divisione, ma prima di prendere una decisione definitiva, vorrei chiederti: sei soddisfatto di questa proposta o ci sono ancora delle possibili modifiche che potremmo apportare per rendere la divisione ancora più equa? 
	  } 
	   } 
	   } 
	 & & \\ 
 

    \theutterance \stepcounter{utterance}  

    & & \multicolumn{4}{p{0.6\linewidth}}{\cellcolor[rgb]{0.9,0.9,0.9}{%
	\makecell[{{p{\linewidth}}}]{% 
	  \tt {\tiny [GM$\rangle$B]}  
	 La tua controproposta è interessante e sembra più equilibrata. Io avrei la corda e due ditali, mentre tu avresti la lampada e un ditale. Questo modo di dividere gli oggetti sembra più giusto e potrebbe funzionare. Sono disposto ad accettare questa divisione, ma prima di prendere una decisione definitiva, vorrei chiederti: sei soddisfatto di questa proposta o ci sono ancora delle possibili modifiche che potremmo apportare per rendere la divisione ancora più equa? 
	  } 
	   } 
	   } 
	 & \\ 
 

    \theutterance \stepcounter{utterance}  

    & & & \multicolumn{4}{p{0.6\linewidth}}{\cellcolor[rgb]{1,0.85,0.72}{%
	\makecell[{{p{\linewidth}}}]{% 
	  \tt {\tiny [GM$\langle$B]}  
	 Sono abbastanza soddisfatto della proposta attuale, ma sono sempre aperto a ulteriori discussioni per trovare la soluzione ottimale. Tuttavia, considerando gli oggetti a disposizione e le nostre esigenze, credo che la divisione proposta sia abbastanza equilibrata. La lampada sembra essere l'oggetto più prezioso, quindi averla io con un ditale potrebbe essere un buon compromesso, lasciandoti la corda e due ditali. Non vedo altre modifiche significative che potrebbero migliorare la divisione senza svantaggiare drasticamente uno di noi. Quindi, a meno che tu non abbia un'ultima proposta o modifica da suggerire, penso che potremmo procedere con l'accordo attuale. [Proposta: 1 lampada, 1 ditale] 
	  } 
	   } 
	   } 
	  \\ 
 

    \theutterance \stepcounter{utterance}  

    & & & \multicolumn{2}{p{0.3\linewidth}}{\cellcolor[rgb]{0.95,0.95,0.95}{%
	\makecell[{{p{\linewidth}}}]{% 
	  \tt {\tiny [GM$|$GM]}  
	 proposal 
	  } 
	   } 
	   } 
	 & & \\ 
 

    \theutterance \stepcounter{utterance}  

    & & & \multicolumn{2}{p{0.3\linewidth}}{\cellcolor[rgb]{0.95,0.95,0.95}{%
	\makecell[{{p{\linewidth}}}]{% 
	  \tt {\tiny [GM$|$GM]}  
	 [1, 1, 0] 
	  } 
	   } 
	   } 
	 & & \\ 
 

    \theutterance \stepcounter{utterance}  

    & & \multicolumn{4}{p{0.6\linewidth}}{\cellcolor[rgb]{0.9,0.9,0.9}{%
	\makecell[{{p{\linewidth}}}]{% 
	  \tt {\tiny [A$\langle$GM]}  
	 L'altro giocatore ha presentato una proposta segreta. Ora tocca a Lei presentare una sua proposta.\\ \tt  
	  } 
	   } 
	   } 
	 & \\ 
 

    \theutterance \stepcounter{utterance}  

    & \multicolumn{4}{p{0.6\linewidth}}{\cellcolor[rgb]{0.8,1,0.9}{%
	\makecell[{{p{\linewidth}}}]{% 
	  \tt {\tiny [A$\rangle$GM]}  
	 [Proposta: 1 lampada, 2 ditali] 
	  } 
	   } 
	   } 
	 & & \\ 
 

    \theutterance \stepcounter{utterance}  

    & & & \multicolumn{2}{p{0.3\linewidth}}{\cellcolor[rgb]{0.95,0.95,0.95}{%
	\makecell[{{p{\linewidth}}}]{% 
	  \tt {\tiny [GM$|$GM]}  
	 proposal 
	  } 
	   } 
	   } 
	 & & \\ 
 

    \theutterance \stepcounter{utterance}  

    & & & \multicolumn{2}{p{0.3\linewidth}}{\cellcolor[rgb]{0.95,0.95,0.95}{%
	\makecell[{{p{\linewidth}}}]{% 
	  \tt {\tiny [GM$|$GM]}  
	 [1, 2, 0] 
	  } 
	   } 
	   } 
	 & & \\ 
 

    \theutterance \stepcounter{utterance}  

    & & & \multicolumn{2}{p{0.3\linewidth}}{\cellcolor[rgb]{0.95,0.95,0.95}{%
	\makecell[{{p{\linewidth}}}]{% 
	  \tt {\tiny [GM$|$GM]}  
	 [[1, 2, 0], [1, 1, 0]] 
	  } 
	   } 
	   } 
	 & & \\ 
 

\end{supertabular}
}

\end{document}
